%% kapa: exam=1 round=19 year=2008 period=1 form=A subject=appraisal_law subject_ko=부동산관계법규 seq=15 number=55
%% source: Q-Net 감정평가사 기출문제 (cst00309) artlSeq=1001962 file=extracted/1008/1001962/2008년 제19회 감정평가사 1차시험 기출문제/1교시.hwp
%% conversion: kordoc 4.24.51 -> tools/gen1.py (auto); review=proofread-sonnet; flags=none
\begin{question}{55}
국토의 계획 및 이용에 관한 법령상 토지거래허가처분에 따른 토지이용의무의 해태에 대하여 부과하는 이행강제금에 관한 설명으로 옳지 않은 것은?
\begin{choices}
  \item 이행강제금 부과의 기준이 되는 토지취득가액은 원칙적으로 실거래가로 한다.
  \item 시장·군수 또는 구청장은 최초의 이행명령이 있은 날을 기준으로 하여 1년에 2회씩 당해 이행명령이 이행될 때까지 반복하여 이행강제금을 부과할 수 있다.
  \item 시장·군수 또는 구청장은 토지거래허가에 따른 이용의무기간이 경과한 후에는 이행강제금을 부과할 수 없다.
  \item 이미 부과된 이행강제금은 그 후 이행명령을 이행하더라도 이를 징수하여야 한다.
  \item 이행강제금의 부과처분에 불복이 있는 자는 부과처분의 고지를 받은 날부터 30일 이내에 시장·군수 또는 구청장에게 이의를 제기할 수 있다.
\end{choices}
\end{question}
\answer{2}
